Fix an integer $k\ge 2$.
For each $n\ge 1$, define $f_k(n)$ to be the smallest integer $M$ such that:

For every set $P$ of $n$ points in $\mathbb{R}^2$ with the property that no $k+1$ points of $P$ are collinear, the number of distinct lines that contain at least $k$ points of $P$ is at most $M$.

Equivalently,
\[
 f_k(n)=\max\{\#\{\ell: |\ell\cap P|\ge k\}:\ |P|=n,\ P\subset\mathbb{R}^2,\ \text{no }k+1\text{ points collinear}\}.
\]
Under the ``no $k+1$ collinear'' assumption, ``at least $k$'' means \emph{exactly $k$}.

Question: for fixed $k\ge 4$, is $f_k(n)=o(n^2)$ as $n\to\infty$?
